\documentclass{article}
\usepackage[T1]{fontenc}
\usepackage{lmodern}
\usepackage{graphicx}
\usepackage{amsmath,amssymb}
\usepackage[a4paper,margin=2cm]{geometry}
\usepackage[absolute,overlay]{textpos}
\setlength{\TPHorizModule}{1mm}
\setlength{\TPVertModule}{1mm}
\usepackage[indonesian]{babel}
\usepackage{hyperref}
\setlength{\emergencystretch}{2em}
% unit: Halaman 54

\begin{document}

% === Halaman 54 (python/native) ===

Galois Field GF(2m)

• Disebut juga medan berhingga biner.  

• GF(2m) atau F2

m adalah ruang vektor berdimensi m pada GF(2).  Setiap elemen di dalam 

GF(2m) adalah integer dalam representasi biner sepanjang maksimal m bit. 

    Contoh: GF(24), GF(28), dst 

• String biner bm-1bm-2 … b1b0,  bi  \{0,1\}, dapat dinyatakan sebagai polinom

            bm-1xm-1 + bm-2xm-2 + … + b1x + b0

• Jadi, setiap a  GF(2m) dapat dinyatakan sebagai 


 bm-1xm-1 + bm-2xm-2 + … + b1x + b0 

• Contoh:  1101  di dalam GF(24) dapat dinyatakan sebagai polinom x3 + x2 + 1 

                    01010111 di dalam GF(28) dapat dinyatakan sebagai polinom x6 + x4 + x2 + x + 1

\end{document}
